% Class-compatibility preamble retained for a reader excerpt of
%   Sahar Diskin, Anna Geisler, "Perfect Matching in Product Graphs and in their
%   Random Subgraphs", arXiv:2404.14020v4, CC BY 4.0.
%
% The source pins its own publisher class (scrartcl); the delivery wrapper accepts
% only the classes article and amsart, so this file supplies the amsart baseline and
% then copies the author's own declaration block mechanically and unchanged
% (pinned-original-source.tex, source lines 30-57: the author's newtheorem
% declarations and the author's macro declarations).  The author's own package
% list of source lines 4-21 is reused only for the entries the retained text needs
% (enumitem, amsmath, amssymb, amsthm, mathtools, xcolor, geometry) and the
% cross-reference packages already needed by the author's own \Cref/\eqref usage
% (hyperref, cleveref); tikz, caption, subcaption, float, graphicx, wrapfig and
% abstract serve figures and the title block that this excerpt does not carry.
% No mathematical macro, statement, symbol or assumption is introduced, changed
% or removed by this file.
\documentclass[10pt,a4paper]{amsart}
\usepackage[shortlabels]{enumitem}
\usepackage{amsmath}
\usepackage{amssymb}
\usepackage{amsthm}
\usepackage{mathtools}
\usepackage{xcolor}
\usepackage[margin=1in]{geometry}
\usepackage{hyperref}
\usepackage{cleveref}
\hypersetup{
  colorlinks=true,
  linkcolor=blue,
  filecolor=magenta,
  urlcolor=cyan
}

\newtheorem{theorem}{Theorem}
\newtheorem{thm}{Theorem}[section]
\newtheorem{corollary}[thm]{Corollary}
\newtheorem{lemma}[thm]{Lemma}
\newtheorem{proposition}[thm]{Proposition}
\newtheorem{claim}[thm]{Claim}
\newtheorem{remark}[thm]{Remark}
\newtheorem{question}[thm]{Question}

\newcommand{\whp}{whp}
\newcommand{\defi}{\operatorname{def}}
\newcommand{\odd}{\operatorname{odd}}
\newcommand{\iso}{\operatorname{iso}}
\newcommand{\cN}{\mathcal{N}}
\newcommand{\cO}{\mathcal{O}}
\newcommand{\cE}{\mathcal{E}}
\newcommand{\cK}{\mathcal{K}}
\newcommand{\cP}{\mathcal{P}}
\newcommand{\cQ}{\mathcal{Q}}

\newcommand{\footremember}[2]{%
  \footnote{#2}%
  \newcounter{#1}%
  \setcounter{#1}{\value{footnote}}%
}

\newcommand{\ag}[1]{\textcolor{purple}{(Anna: #1)}}
\newcommand{\sd}[1]{\textcolor{red}{(Sahar: #1)}}
\crefname{originallemma}{original lemma}{original lemmas}
\Crefname{originallemma}{Original lemma}{Original lemmas}
\makeatletter
\newlabel{lem:three_edges}{{6.2}{source}{}{sourcepointer-1}{}}
\newlabel{lem:three_edges@cref}{{[originallemma][0][]6.2}{[1][0][]source}{}{}{}}
\makeatother
\crefname{originallemma}{original lemma}{original lemmas}
\Crefname{originallemma}{Original lemma}{Original lemmas}
\makeatletter
\newlabel{lem:small-boundary-sets}{{6.3}{source}{}{sourcepointer-2}{}}
\newlabel{lem:small-boundary-sets@cref}{{[originallemma][0][]6.3}{[1][0][]source}{}{}{}}
\makeatother
\crefname{originalequation}{original equation}{original equations}
\Crefname{originalequation}{Original equation}{Original equations}
\makeatletter
\newlabel{eq:mu-range}{{4.2}{source}{}{sourcepointer-3}{}}
\newlabel{eq:mu-range@cref}{{[originalequation][0][]4.2}{[1][0][]source}{}{}{}}
\makeatother
\crefname{originalequation}{original equation}{original equations}
\Crefname{originalequation}{Original equation}{Original equations}
\makeatletter
\newlabel{eq:Gallai-Edmonds}{{3.1}{source}{}{sourcepointer-4}{}}
\newlabel{eq:Gallai-Edmonds@cref}{{[originalequation][0][]3.1}{[1][0][]source}{}{}{}}
\makeatother
\crefname{originallemma}{original lemma}{original lemmas}
\Crefname{originallemma}{Original lemma}{Original lemmas}
\makeatletter
\newlabel{lem:no-small-separator}{{6.5}{source}{}{sourcepointer-5}{}}
\newlabel{lem:no-small-separator@cref}{{[originallemma][0][]6.5}{[1][0][]source}{}{}{}}
\makeatother
\crefname{originallemma}{original lemma}{original lemmas}
\Crefname{originallemma}{Original lemma}{Original lemmas}
\makeatletter
\newlabel{lem:small-components-count}{{6.6}{source}{}{sourcepointer-6}{}}
\newlabel{lem:small-components-count@cref}{{[originallemma][0][]6.6}{[1][0][]source}{}{}{}}
\makeatother
\title{With high probability the random subgraph admits no small Tutte obstruction in the whole window of edge probabilities}
\author{Sahar Diskin, Anna Geisler}
\date{arxiv:2404.14020v4; CC BY 4.0}
\begin{document}
\maketitle
\noindent\textit{Editorial note.} All delivered mathematical text is the author's own, copied from the pinned author source sources/combinatorics-fl04-vspace-1-5cm-perfect-matching-in-p/pinned-original-source.tex; no statement or proof step was rewritten, reordered or removed.\par\medskip
\noindent\textit{Editorial note.} Preamble provenance: the source pins the publisher class scrartcl, which the delivery wrapper does not accept, so this card uses sources/combinatorics-fl04-vspace-1-5cm-perfect-matching-in-p/retained-preamble.tex, which supplies the amsart baseline and then carries the author's own newtheorem and macro declarations (pinned source lines 30-57) verbatim. No macro was invented and no mathematical symbol was changed.\par\medskip
\noindent\textit{Editorial note.} Bibliography provenance: the source carries no thebibliography environment and ends with bibliographystyle abbrv plus bibliography perc (author source lines 1691-1692), so the bibliography is resolved with BibTeX from the identical copy staged with the same arXiv version under sources/product-graph-matchings-2404.14020v4/perc.bib; no entry was edited and the bibliography style is the author's own abbrv.\par\medskip
\noindent\textit{Editorial note.} Two labels the counted proof invokes, lem:three\_edges and lem:small-boundary-sets, and the displayed range for p are navigation pointers only; the content they name is declared as named standard prerequisites and is not reproduced here. The same is true of the Gallai-Edmonds relation eq:Gallai-Edmonds, which enters only through the decomposition the surrounding subsection sets up.\par\medskip
\noindent\textit{Editorial note.} The counted statement is delivered together with the Tutte-violation display and the definition of u\_0 as context spans, both quoted verbatim from the author's own text.\par\medskip
\noindent\textit{Allowed standard prerequisites (editorial).}\begin{itemize}
\item The Tutte-violation relation and the threshold u\_0. Stated by the author immediately before the two lemmas (source lines 656-668), delivered verbatim as context spans of this card. Fixes the property that the obstruction set must satisfy and the size threshold below which the lemma excludes it.
\item The minimality lemma for obstructions and the counting lemmas for sets with a small edge boundary. The author's lemmas at source lines 703-705 and 714-716, whose proofs are delivered by this lane's cards for those statements. Supplies the two ingredients whose combination gives the size bound for the obstruction.
\end{itemize}\medskip
\ifdefined\hypertarget\hypertarget{sourcepointer-1}{}\fi
\noindent\textit{Source reference.} Original lemma 6.2: sources/combinatorics-fl04-vspace-1-5cm-perfect-matching-in-p/pinned-original-source.tex, source lines 703--703. Navigation pointer only: the counted proof uses this minimality lemma in its first step; its own statement and proof are delivered by this lane's separate card for that statement and are not part of this card.\par\medskip
\ifdefined\hypertarget\hypertarget{sourcepointer-2}{}\fi
\noindent\textit{Source reference.} Original lemma 6.3: sources/combinatorics-fl04-vspace-1-5cm-perfect-matching-in-p/pinned-original-source.tex, source lines 714--714. Navigation pointer only: the counted proof uses this counting lemma in its second step; its own statement and proof are delivered by this lane's separate card and are not part of this card.\par\medskip
\ifdefined\hypertarget\hypertarget{sourcepointer-3}{}\fi
\noindent\textit{Source reference.} Original equation 4.2: sources/combinatorics-fl04-vspace-1-5cm-perfect-matching-in-p/pinned-original-source.tex, source lines 479--479. Navigation pointer only: the range of p displayed there is what makes the estimate uniform over the whole window; the display is not part of this card.\par\medskip
\ifdefined\hypertarget\hypertarget{sourcepointer-4}{}\fi
\noindent\textit{Source reference.} Original equation 3.1: sources/combinatorics-fl04-vspace-1-5cm-perfect-matching-in-p/pinned-original-source.tex, source lines 266--266. Navigation pointer only: the counted proof consumes the Gallai-Edmonds decomposition that the surrounding subsection sets up; the relation itself is not part of this card.\par\medskip
\ifdefined\hypertarget\hypertarget{sourcepointer-5}{}\fi
\noindent\textit{Source reference.} Original lemma 6.5: sources/combinatorics-fl04-vspace-1-5cm-perfect-matching-in-p/pinned-original-source.tex, source lines 772--772. Navigation pointer only: the counted proof consumes the statement or estimate labelled here; its content is declared as a named standard prerequisite and the labelled text is not reproduced in this card.\par\medskip
\ifdefined\hypertarget\hypertarget{sourcepointer-6}{}\fi
\noindent\textit{Source reference.} Original lemma 6.6: sources/combinatorics-fl04-vspace-1-5cm-perfect-matching-in-p/pinned-original-source.tex, source lines 842--842. Navigation pointer only: the counted proof consumes the statement or estimate labelled here; its content is declared as a named standard prerequisite and the labelled text is not reproduced in this card.\par\medskip
\begin{lemma}\label{lem:no-small-tutte}
Whp, simultaneously for every $p\in[p_-,p_+]$ with $\delta(G_p)\ge1$, the graph $G_p$ has no set $U$ with $1\le|U|\le u_0$ satisfying \eqref{eq:tutte-violation}.
\end{lemma}

\begin{align} \label{eq:tutte-violation}
    \odd(G_p - U) \geq |U|+2
\end{align}

\[
 u_0\coloneqq \frac{n}{d^{M^5/p_-}}.
\]


\begin{proof}[Proof of \Cref{lem:no-small-tutte}]
Recall that whp every $G_p$ under consideration is connected and that $G$ has even order. In particular, the empty set does not satisfy \eqref{eq:tutte-violation}.
Choose a set $U$ satisfying \eqref{eq:tutte-violation} of minimum cardinality, put $u=|U|$, and let $K_1,\ldots,K_m$ be the odd components of $G_p-U$.
By \Cref{lem:three_edges} every $v\in U$ has neighbours in at least three of these components. 
Every $K_i$ which is an isolated vertex in $G_p-U$ has a $G_p$-neighbour in $U$, because $\delta(G_p)\ge1$.

Suppose that $u \leq u_0$.  
Whp, at most one $K_i$ has order larger than $u_1$ by \Cref{lem:no-small-separator}. Let $B$ be the union of the remaining odd components. Note that there are at least $u+1$ of them and every $v\in U$ has neighbours in at least two components of $B$.
If $|B| < u_1$, then $U$ and $B$ satisfy the assumptions of \Cref{lem:small-components-count}, a contradiction.
If $|B|\geq u_1$, then a greedy union of its components has size between $u_1$ and $2u_1<n/2$. In this case $U$ together with this union of components satisfy the assumptions of \Cref{lem:no-small-separator}, a contradiction.
\end{proof}





\bibliographystyle{plain}
\bibliography{perc}
\end{document}
